On plonge la cellule d'un conductimètre ($S=\qty{1,0}{\square\cm}~$;
$\ell=\qty{1,0}{\cm}$) dans une solution d'acide méthanoïque \ce{HCOOH} de
concentration $C=\qty{5,0e-3}{\M}$. On lit une conductance
$G=\qty{3,30e-4}{\siemens}$.

\begin{enumerate}
    \item Écrire l'équation de la réaction de l'acide méthanoïque avec l'eau et
          dresser le tableau d'avancement.
    \item Déterminer la valeur de la conductivité $\sigma$.
    \item Calculer les concentrations molaires des ions produits par cette
          réaction.
    \item Déterminer le taux d'avancement final $\tau$ de la réaction de l'acide
          méthanoïque avec l'eau.
    \item Calculer la constante d'équilibre $\mathrm{K}$ associée à l'équation
          de cette réaction.
\end{enumerate}

\begin{donnees}
    \begin{itemize}
        \item Conductivités molaires ioniques (en
              $\unit{\siemens\square\meter\per\mole}$)~: \par
              $\lambda_{\ce{H3O+}}=\num{3,50e-2}$ \;; \;
              $\lambda_{\ce{HCOO-}}=\num{5,50e-3}$
    \end{itemize}
\end{donnees}

